\section{Introduction and scope}
\label{sec:introduction}

The factor-ray transform separates two different kinds of averaging.
Continuous Mellin averaging separates distinct reduced ratios $a/b$,
whereas arithmetic averaging of the shift $h$ controls the prime targets
that remain coherent on a single ray.  TPC-11 established the first
separation at fixed shift and evaluated the resulting ray energy after a
smooth shift average of length $H\geq X$ \citep{WangTPC11}.  Its explicit
next question was whether one can shorten the shift average without
pretending that a fixed shift has been reached.

The recent large-values theorem of \citet{GuthMaynard2026} changes the
answer.  Among its consequences is a prime number theorem in intervals of
length $X^{2/15+\eps}$ for all but a stretched-exponentially sparse set of
integer starting points.  The exponent $2/15$ applies to almost all
intervals; it is logically different from the exponent $17/30$ in their
uniform short-interval theorem.  The exceptional-set framework and its
current quantitative refinements are also discussed by
\citet{GafniTao2026}.  We use the almost-all result because the ray energy
already sums over the product center $r=mn$.

That last sentence contains the main technical issue.  A set that is
sparse among all integers could, in principle, carry most of a structured
divisor weight.  The relevant weight here is not bounded:
\[
 A_D(r)=|W(r/X)|^2
 \sum_{\substack{n\mid r\\n>D}}(\Lambda(n)-1)^2.
\]
We close this gap by proving a global second-moment bound for $A_D$ and
then transferring the exceptional-set estimate by Cauchy--Schwarz.  This
weighted transfer, rather than a pointwise short-interval theorem, is what
allows the shift scale to fall below the uniform threshold.

\subsection{Main results}

Throughout the paper
\[
 u(n)=\Lambda(n)-1,
 \qquad
 \eta(q)=\Lambda(q)-1,
\]
and $W\in C_c^1((0,\infty);\R)$ has support in a fixed interval
$[\alpha,\beta]\subset(0,\infty)$.  We write
\begin{equation}
 \cT_{h,D}(X;t)=
 \sum_{\substack{n>D\\m\geq1}}
 u(n)\Lambda(mn+h)W(mn/X)(n/m)^{it},
 \label{eq:intro-raw-transform}
\end{equation}
\begin{equation}
 \cM_D(X;t)=
 \sum_{\substack{n>D\\m\geq1}}
 u(n)W(mn/X)(n/m)^{it},
 \qquad
 \cR_{h,D}=\cT_{h,D}-\cM_D.
 \label{eq:intro-centered-transform}
\end{equation}
The factorization $n=ak$, $m=bk$, $(a,b)=1$, gives
\begin{equation}
 \cR_{h,D}(X;t)=
 \sum_{(a,b)=1}c_{a,b}(h;X,D)(a/b)^{it},
 \label{eq:intro-ray-compression}
\end{equation}
where
\begin{equation}
 c_{a,b}(h;X,D)=
 \sum_{\substack{k\geq1\\ak>D}}
 u(ak)\eta(abk^2+h)W(abk^2/X).
 \label{eq:intro-ray-coefficient}
\end{equation}
The centered ray energy is
\begin{equation}
 \cE_{h,D}(X)=\sum_{(a,b)=1}|c_{a,b}(h;X,D)|^2.
 \label{eq:intro-ray-energy}
\end{equation}

Our first theorem transfers the almost-all short-interval prime number
theorem through both the signed first-moment coefficients and the positive
diagonal source weights.  In particular, for every $A>0$,
\begin{equation}
 \frac1H\sum_hG(h/H)\cT_{h,D}(X;t)
 =I_0(G)\cM_D(X;t)+O_A\bigl(X(\log X)^{-A}\bigr)
 \label{eq:intro-first-moment}
\end{equation}
uniformly in $t\in\R$ throughout the growing parameter range of the
energy theorem below.

The main energy theorem is the following.  Put
\begin{equation}
 I_0(G)=\int_0^\infty G(v)\,dv,
 \qquad
 I_2(W)=\int_0^\infty |W(v)|^2\,dv.
 \label{eq:intro-moments}
\end{equation}
For fixed $0<\eps<13/15$ and $0<\eta<1$, uniformly in
\[
 X^{2/15+\eps}\leq H\leq X,
 \qquad
 1\leq D\leq X^{1-\eta},
\]
we prove
\begin{equation}
 \begin{aligned}
 \sum_hG(h/H)\cE_{h,D}(X)
 ={}&\frac12HX I_0(G)I_2(W)\log X\\
 &\times\bigl((\log X)^2-(\log D)^2\bigr)\\
 &+O_{\eps,\eta,W,G}\!\left(
 HX(\log X)^2\log\log(3X)\right).
 \end{aligned}
 \label{eq:intro-main-energy}
\end{equation}
The raw ray energy has the same asymptotic.  In this at-most-linear shift
range the correct target logarithm is $\log X$, not $\log H$.

The source diagonal is of order $X(\log X)^2$.  On a nonexceptional
center, the target square contributes $H\log X$.  A two-linear-forms
Selberg upper sieve bounds each same-ray off-diagonal by
$H\log\log X$, after which the source off-diagonal ledger from TPC-11 is
$O(X\log^2X)$.  Thus the extra logarithm on the diagonal remains the sole
source of the asymptotic.  The proof does not obtain cancellation in an
individual fixed-shift quadratic correlation.

The exact factor-ray mean-square identities then convert
\eqref{eq:intro-main-energy} into hard- and Fej\'er-band results over the
same short shift window.  We also prove a parity-conditioned version over
even shifts.  Its main coefficient is $1/4$ rather than $1/2$, and it
does not select any prescribed even shift.

\subsection{The fixed-shift boundary}

Shrinking a growing average and setting $H=1$ are not interchangeable
operations.  We make that failure concrete in three ways.

\begin{enumerate}[label=\textup{(\roman*)}]
 \item For nonnegative, nonzero $W$ and every fixed $h_0\geq2$, there is
 an infinite family of
 primitive rays on which the actual arithmetic summands in
 \eqref{eq:intro-ray-coefficient} have one sign and no cancellation at
 all.  This includes balanced ray sectors whenever the cutoff is below
 the square-root scale; a separate unbalanced construction covers the
 full range $D\leq X^{1-\eta}$.

 \item On the locally admissible ray $(a,b,h)=(6,1,1)$, the unknown part
 of the coefficient is exactly
 \[
  \sum_k\Lambda(6k^2+1)W(6k^2/X).
 \]
 Therefore a natural-scale coefficient asymptotic is equivalent to the
 corresponding weighted quadratic-prime asymptotic.  Bateman--Horn
 predicts such an asymptotic \citep{BatemanHorn1962}, but it is not known;
 the almost-prime theorem for $n^2+1$ does not supply it
 \citep{Iwaniec1978}.

 \item The whole one-Mellin signal on one ray is a known exponential
 times the single scalar $\sum_kx_k$.  We give the exact minimax error for
 recovering any other radial functional from that scalar.  This is an
 information statement for coefficient-blind recovery, not a prohibition
 on using additional arithmetic input.
\end{enumerate}

These statements identify both a local obstruction and an admissible open
problem before the endpoint prime-pair layer is reached.  At the endpoint,
the selector in TPC-11 already shows that a sufficiently strong
fixed-shift band approximation would recover a weighted Hardy--Littlewood
functional \citep{HardyLittlewood1923,WangTPC11}.

\subsection{Claim ledger}

\begin{center}
\begin{tabular}{@{}p{0.31\textwidth}p{0.27\textwidth}p{0.31\textwidth}@{}}
\toprule
Statement & Status here & What it does not imply\\
\midrule
Weighted transfer of almost-all short-interval PNT
 & Proved in the displayed growing range
 & No prescribed shift is controlled\\
Short-shift factor-ray energy asymptotic
 & Proved with explicit error
 & No individual ray or target pair is evaluated\\
Even-shift ensemble
 & Proved after parity centering
 & No selection of a particular even shift\\
Fixed-ray cancellation
 & False on an explicit local-obstruction family
 & Does not contradict averaged energy\\
Admissible quadratic ray
 & Exact algebraic reduction
 & The quadratic-prime asymptotic remains open\\
\bottomrule
\end{tabular}
\end{center}

No twin-prime lower bound, fixed-shift Hardy--Littlewood asymptotic,
quadratic-prime theorem, new level of distribution, or breach of the
sieve parity barrier is asserted.
